\chapter{Если что-то пошло не так}
\label{ch:troubleshooting}

Здесь собраны ситуации, которые выглядят как поломка сайта, но объясняются проще. Многие из них --- следствие того, что интерфейс пока не показывает причину происходящего; в таких случаях ниже сказано прямо, что именно означает пустой экран.

\section{В сессии пусто, ничего не появляется}
\label{sec:trouble:empty-session}

\paragraph{Скорее всего, вам не назначен персонаж.}
Экран игрока без персонажа остаётся пустым --- без надписи и без кнопок. Это не сбой: мастер просто ещё не выдал вам героя. Напишите ему.

\section{Бесконечная крутилка при входе в сессию}
\label{sec:trouble:spinner}

Загрузка, которая не заканчивается, означает, что сессия не прислала данные. Причин три:

\begin{enumerate}
\item игра ещё не запущена --- мастер держит всех в лобби;
\item подход уже завершён, и ссылка устарела;
\item у вас нет доступа именно к этой игре.
\end{enumerate}

Таймаута у экрана нет, ждать бесполезно. Обновите страницу; если не помогло --- уточните у мастера, идёт ли игра и та ли это ссылка.

\section{``У тебя нет роли''}
\label{sec:trouble:no-role}

Вы открыли сессию, в которой не участвуете: не мастер и не игрок. Обычно это чужая ссылка или ссылка на другую игру. Попросите мастера добавить вас в лобби --- или дать ссылку наблюдателя, если участвовать не планировалось.

\section{Нажал кнопку --- ничего не произошло}
\label{sec:trouble:connection}

Признак оборвавшейся связи. Экран остаётся заполненным, кнопки нажимаются, но до сервера ничего не доходит и новые события не приходят. Индикатора соединения на экране пока нет, поэтому ориентируйтесь косвенно: если несколько минут никто не двигается и ваши действия не появляются в ленте --- связь потеряна.

Лечится перезагрузкой страницы. Игра при этом не теряется: состояние хранится на сервере, вы вернётесь туда же.

\section{Список лобби пуст}
\label{sec:trouble:empty-lobbies}

``Лобби пока нет'' и ``не удалось получить список'' выглядят одинаково --- пустым списком. Обновите страницу; если пусто и после обновления, а мастер уверяет, что лобби создано --- попросите у него прямую ссылку вида \texttt{/lobby/<id>}.

\section{Пропала кнопка отмены действия}
\label{sec:trouble:cancel-gone}

Бросок уже сделан. После броска игрок отменить действие не может --- результат записан. Отменить может мастер, попросите его.

\section{Не могу создать сценарий}
\label{sec:trouble:no-master}

Раздел открывается, но создавать нечего --- значит, у вас нет прав мастера. Их выдаёт администратор платформы, см.\ раздел \ref{sec:admin:grant-master}. После выдачи нужно перезайти на сайт.

\section{Не работает вход через Telegram в браузере}
\label{sec:trouble:tg-login}

Кнопки ``Войти через Telegram'' на странице входа нет --- такого способа в обычном браузере не существует. Есть два рабочих пути: открыть сайт кнопкой прямо из Telegram (там вход произойдёт сам) либо запросить у бота одноразовую ссылку командой \texttt{/link}.

\section{Наблюдатель видит пустой экран}
\label{sec:trouble:observer-empty}

Либо мастер ещё не открыл ни одной сцены, либо код комнаты введён неверно. Сообщения об ошибке не будет ни в том, ни в другом случае --- сверьте код по символам и спросите мастера, началась ли игра.

\section{Правки локации не сохранились}
\label{sec:trouble:lost-edits}

Изменения локаций, сделанные во время сессии, могут не дойти до базы и пропасть при перезагрузке. Вносите такие правки в редакторе сценария, а не на ходу.

\section{Выкидывает на страницу входа}
\label{sec:trouble:logout}

Истёк срок сессии. Войдите заново --- вас вернёт на ту же страницу. Если выкидывает сразу после входа, повторяющимся кругом, очистите куки сайта и попробуйте ещё раз.
